% GENERATED FILE. Edit canonical lessons, then run npm run generate:books.
% canonical-source-hash: fnv1a64:12b992ca02ca5464
% canonical-lessons: JA-C422-washoku, JA-C422-youshoku, JA-C422-teishoku, JA-C422-sashimi, JA-C422-yakitori

\chapter{Eating Out on the Road}
\label{ch:ja-eating-out-on-the-road}
\hlchaptermodality{422}

\begin{chapteropening}
\textbf{By the end of this chapter:} \emph{I can choose what to eat on a trip.}

Complete the last lesson of chapter 422: I can choose what to eat on a trip.
\end{chapteropening}

\section[\texorpdfstring{\ja{わしょく} (washoku)}{わしょく (washoku)}]{\ja{わしょく} (washoku) — food in the Japanese style}

\label{lesson:JA-C422-washoku}



\begin{quote}
Before the new one: say the Japanese for a price cut, then the Japanese for merchandise.
\end{quote}

\subsection*{You'll want to know: \ja{わしょく}}
\textbf{\ja{わしょく}} — \emph{washoku} — \textquotedblleft{}food in the Japanese style\textquotedblright{}.

\begin{quote}
\textbf{\ja{たびさきで} \ja{わしょくを} \ja{たべました。}} — \emph{tabisaki de washoku o tabemashita} — \textquotedblleft{}I ate Japanese food where I was staying.\textquotedblright{}
\end{quote}

\begin{cousinweb}
\textbf{\ja{わしょく}} (\emph{washoku}) is Sino-Japanese (kango): two Chinese characters read with their on'yomi, readings borrowed from Chinese. They are \emph{wa}, \textquotedblleft{}Japanese style\textquotedblright{} (Middle Chinese \emph{hwa}) and \emph{shoku}, \textquotedblleft{}food\textquotedblright{} (\emph{zyik}) (KANJIDIC2; Middle Chinese in Baxter's transcription).
\end{cousinweb}

\begin{grammarlens}[title={What you've built}]
One more word for getting around.
\end{grammarlens}

\subsection*{Your turn}
\begin{itemize}
\raggedright
  \item \emph{Say it:} \emph{washoku}
  \item \emph{Say it:} \emph{washoku}, once more
  \item \emph{Say it:} \emph{tabisaki de washoku o tabemashita} — I ate Japanese food where I was staying.
  \item \emph{Recall:} say the Japanese for a price cut, then the Japanese for merchandise, then say \emph{washoku} again
\end{itemize}

\begin{tcolorbox}[breakable,colback=teal!4,colframe=teal!35!black,title={Before you move on}]
What does \ja{わしょく} mean? (\textquotedblleft{}Food in the Japanese style\textquotedblright{}.) Say it once more.
\end{tcolorbox}
\section[\texorpdfstring{\ja{ようしょく} (yōshoku)}{ようしょく (yōshoku)}]{\ja{ようしょく} (yōshoku) — food in the Western style}

\label{lesson:JA-C422-youshoku}



\begin{quote}
Before the new one: say the Japanese for merchandise, then the Japanese for food in the Japanese style.
\end{quote}

\subsection*{You'll want to know: \ja{ようしょく}}
\textbf{\ja{ようしょく}} — \emph{yōshoku} — \textquotedblleft{}food in the Western style\textquotedblright{}.

\begin{quote}
\textbf{\ja{あには} \ja{ようしょくより} \ja{わしょくが} \ja{すきです。}} — \emph{ani wa yōshoku yori washoku ga suki desu} — \textquotedblleft{}My brother likes Japanese food better than food in the western style.\textquotedblright{}
\end{quote}

\begin{cousinweb}
\textbf{\ja{ようしょく}} (\emph{yōshoku}) is Sino-Japanese (kango): two Chinese characters read with their on'yomi, readings borrowed from Chinese. They are \emph{yō}, \textquotedblleft{}ocean\textquotedblright{} (Middle Chinese \emph{yang}) and \emph{shoku}, \textquotedblleft{}food\textquotedblright{} (\emph{zyik}) (KANJIDIC2; Middle Chinese in Baxter's transcription).
\end{cousinweb}

\begin{grammarlens}[title={What you've built}]
One more word for getting around.
\end{grammarlens}

\subsection*{Your turn}
\begin{itemize}
\raggedright
  \item \emph{Say it:} \emph{yōshoku}
  \item \emph{Say it:} \emph{yōshoku}, once more
  \item \emph{Say it:} \emph{ani wa yōshoku yori washoku ga suki desu} — My brother likes Japanese food better than food in the western style.
  \item \emph{Recall:} say the Japanese for merchandise, then the Japanese for food in the Japanese style, then say \emph{yōshoku} again
\end{itemize}

\begin{tcolorbox}[breakable,colback=teal!4,colframe=teal!35!black,title={Before you move on}]
What does \ja{ようしょく} mean? (\textquotedblleft{}Food in the Western style\textquotedblright{}.) Say it once more.
\end{tcolorbox}
\section[\texorpdfstring{\ja{ていしょく} (teishoku)}{ていしょく (teishoku)}]{\ja{ていしょく} (teishoku) — a set meal}

\label{lesson:JA-C422-teishoku}



\begin{quote}
Before the new one: say the Japanese for food in the Japanese style, then the Japanese for food in the Western style.
\end{quote}

\subsection*{You'll want to know: \ja{ていしょく}}
\textbf{\ja{ていしょく}} — \emph{teishoku} — \textquotedblleft{}a set meal\textquotedblright{}.

\begin{quote}
\textbf{\ja{しょくどうで} \ja{さかなの} \ja{ていしょくを} \ja{たのみました。}} — \emph{shokudō de sakana no teishoku o tanomimashita} — \textquotedblleft{}I ordered the fish set meal at the canteen.\textquotedblright{}
\end{quote}

\begin{cousinweb}
\textbf{\ja{ていしょく}} (\emph{teishoku}) is Sino-Japanese (kango): two Chinese characters read with their on'yomi, readings borrowed from Chinese. They are \emph{tei}, \textquotedblleft{}fix\textquotedblright{} (Middle Chinese \emph{dengH}) and \emph{shoku}, \textquotedblleft{}food\textquotedblright{} (\emph{zyik}) (KANJIDIC2; Middle Chinese in Baxter's transcription).
\end{cousinweb}

\begin{grammarlens}[title={What you've built}]
One more word for getting around.
\end{grammarlens}

\subsection*{Your turn}
\begin{itemize}
\raggedright
  \item \emph{Say it:} \emph{teishoku}
  \item \emph{Say it:} \emph{teishoku}, once more
  \item \emph{Say it:} \emph{shokudō de sakana no teishoku o tanomimashita} — I ordered the fish set meal at the canteen.
  \item \emph{Recall:} say the Japanese for food in the Japanese style, then the Japanese for food in the Western style, then say \emph{teishoku} again
\end{itemize}

\begin{tcolorbox}[breakable,colback=teal!4,colframe=teal!35!black,title={Before you move on}]
What does \ja{ていしょく} mean? (\textquotedblleft{}A set meal\textquotedblright{}.) Say it once more.
\end{tcolorbox}
\section[\texorpdfstring{\ja{さしみ} (sashimi)}{さしみ (sashimi)}]{\ja{さしみ} (sashimi) — sliced raw fish}

\label{lesson:JA-C422-sashimi}



\begin{quote}
Before the new one: say the Japanese for food in the Western style, then the Japanese for a set meal.
\end{quote}

\subsection*{You'll want to know: \ja{さしみ}}
\textbf{\ja{さしみ}} — \emph{sashimi} — \textquotedblleft{}sliced raw fish\textquotedblright{}.

\begin{quote}
\textbf{\ja{みなとで} \ja{さしみを} \ja{たべました。}} — \emph{minato de sashimi o tabemashita} — \textquotedblleft{}We ate sliced raw fish at the harbour.\textquotedblright{}
\end{quote}

\begin{cousinweb}
\textbf{\ja{さしみ}} (\emph{sashimi}) is a native Japanese word (wago), inherited rather than borrowed from Chinese, built from \emph{sashi}, \textquotedblleft{}piercing\textquotedblright{} (from \emph{sasu}) and \emph{mi}, \textquotedblleft{}flesh\textquotedblright{}; its kanji were attached later as kun'yomi (JMdict).
\end{cousinweb}

\begin{grammarlens}[title={What you've built}]
One more word for getting around.
\end{grammarlens}

\subsection*{Your turn}
\begin{itemize}
\raggedright
  \item \emph{Say it:} \emph{sashimi}
  \item \emph{Say it:} \emph{sashimi}, once more
  \item \emph{Say it:} \emph{minato de sashimi o tabemashita} — We ate sliced raw fish at the harbour.
  \item \emph{Recall:} say the Japanese for food in the Western style, then the Japanese for a set meal, then say \emph{sashimi} again
\end{itemize}

\begin{tcolorbox}[breakable,colback=teal!4,colframe=teal!35!black,title={Before you move on}]
What does \ja{さしみ} mean? (\textquotedblleft{}Sliced raw fish\textquotedblright{}.) Say it once more.
\end{tcolorbox}
\section[\texorpdfstring{\ja{やきとり} (yakitori)}{やきとり (yakitori)}]{\ja{やきとり} (yakitori) — grilled chicken on skewers}

\label{lesson:JA-C422-yakitori}



\begin{quote}
Before the new one: say the Japanese for a set meal, then the Japanese for sliced raw fish.
\end{quote}

\subsection*{You'll want to know: \ja{やきとり}}
\textbf{\ja{やきとり}} — \emph{yakitori} — \textquotedblleft{}grilled chicken on skewers\textquotedblright{}.

\begin{quote}
\textbf{\ja{えきまえの} \ja{みせで} \ja{やきとりを} \ja{たべました。}} — \emph{ekimae no mise de yakitori o tabemashita} — \textquotedblleft{}We had grilled chicken at a shop in front of the station.\textquotedblright{}
\end{quote}

\begin{cousinweb}
\textbf{\ja{やきとり}} (\emph{yakitori}) is a native Japanese word (wago), inherited rather than borrowed from Chinese, built from \emph{yaki}, \textquotedblleft{}grilling\textquotedblright{} (from \emph{yaku}) and \emph{tori}, \textquotedblleft{}bird\textquotedblright{}; its kanji were attached later as kun'yomi (JMdict).
\end{cousinweb}

\begin{grammarlens}[title={What you've built}]
One more word for getting around.
\end{grammarlens}

\subsection*{Your turn}
\begin{itemize}
\raggedright
  \item \emph{Say it:} \emph{yakitori}
  \item \emph{Say it:} \emph{yakitori}, once more
  \item \emph{Say it:} \emph{ekimae no mise de yakitori o tabemashita} — We had grilled chicken at a shop in front of the station.
  \item \emph{Recall:} say the Japanese for a set meal, then the Japanese for sliced raw fish, then say \emph{yakitori} again
\end{itemize}

\begin{tcolorbox}[breakable,colback=teal!4,colframe=teal!35!black,title={Before you move on}]
What does \ja{やきとり} mean? (\textquotedblleft{}Grilled chicken on skewers\textquotedblright{}.) Say it once more.
\end{tcolorbox}
